% DOCUMENT CLASS SETUP
\documentclass[12pt, a4paper]{article}

% DOCUMENT PACKAGES
\usepackage[a4paper, margin=1in]{geometry}
\usepackage{titling}
\usepackage{enumitem}
\usepackage{amsmath}
\usepackage{amssymb}
\usepackage{mathtools}
\usepackage{graphicx}
\usepackage{xcolor}
\usepackage{listings}
\usepackage{hyperref}
\usepackage{csquotes}

% FONT AND ENCODING SETUP
\usepackage{fontspec}
\usepackage{polyglossia}
\usepackage[Thai]{ucharclasses}
\setdefaultlanguage{english}
\setotherlanguage{thai}
\setTransitionsFor{Thai}{\thaifont}{\normalfont}
\newfontfamily{\thaifont}{THSarabun.ttf}[
    Path=./,
    BoldFont=THSarabunBold.ttf,
    ItalicFont=THSarabunItalic.ttf,
    BoldItalicFont=THSarabunBoldItalic.ttf,
    Scale=1.23
]
\XeTeXlinebreaklocale "th"
\XeTeXlinebreakskip = 0pt plus 1pt

% PROFESSIONAL CODE BLOCK
\usepackage{tcolorbox}
\tcbuselibrary{minted, skins, breakable}

% DEFINE GITHUB DARK THEME COLORS %
\definecolor{ghback}{HTML}{0d1117}    % Dark background
\definecolor{ghtext}{HTML}{c9d1d9}    % Light gray text
\definecolor{ghborder}{HTML}{30363d}  % GitHub Border color

% PROFESSIONAL CODE BLOCK CONFIGURATION (MINTED VERSION) %
\renewcommand{\theFancyVerbLine}{\textcolor{white!60!gray}{\scriptsize\arabic{FancyVerbLine}}}
\newtcblisting{codeblock}[2][]{
    enhanced,
    breakable,
    colback=ghback,
    colframe=ghborder,
    coltitle=ghtext,
    colupper=ghtext,
    fonttitle=\bfseries\small,
    arc=3mm,
    boxrule=0.5mm,
    % PADDING SETTINGS
    left=25pt,
    right=10pt,
    top=5pt,
    bottom=5pt,
    % MINTED SETTINGS
    listing engine=minted,
    listing only,
    minted language=python,
    minted options={
        linenos,
        numbersep=5pt,
        fontsize=\footnotesize,
        breaklines=true,
        breakanywhere=true,
        autogobble,
        style=monokai,
    },
    title={#2},
    #1
}

% DOCUMENT TITLE SETUP %
\title {
	\textbf{2110452 HIGH PERFORMANCE ARCHITECTURE} \\
    \vspace{1em}
	\Large{\textbf{Homework 2: Loop Transformation}} \\
    \vspace{0.5em}
    \large{Computer Engineering, Chulalongkorn Univerity}
}
\author{Worralop Srichainont}
\date{\today}
\setlength{\droptitle}{15em}

% DOCUMENT CONTENTS
\begin{document}

    \maketitle
    \thispagestyle{empty}
    \clearpage

    \tableofcontents
    \thispagestyle{empty}
    \clearpage

    \pagenumbering{arabic}

    \section{Problem I}
        Try code the matrix (e.g. $8 \times 8$) multiplication in a straight forward algorithm.
        Compile it with various optimizations.
        See if the compiler is able to perform loop interchange.

        \subsection{Matrix Multiplication Code}
            \noindent Given the $N \times N$ matrix multiplication below
            $$
                \begin{bmatrix}
                    c_{11} & c_{12} & \dots & c_{1N} \\
                    c_{21} & c_{22} & \dots & c_{2N} \\
                    \vdots & \vdots & \ddots & \vdots \\
                    c_{N1} & c_{N2} & \dots & c_{NN} \\
                \end{bmatrix}
                =
                \begin{bmatrix}
                    a_{11} & a_{12} & \dots & a_{1N} \\
                    a_{21} & a_{22} & \dots & a_{2N} \\
                    \vdots & \vdots & \ddots & \vdots \\
                    a_{N1} & a_{N2} & \dots & a_{NN} \\
                \end{bmatrix}
                \begin{bmatrix}
                    b_{11} & b_{12} & \dots & b_{1N} \\
                    b_{21} & b_{22} & \dots & b_{2N} \\
                    \vdots & \vdots & \ddots & \vdots \\
                    b_{N1} & b_{N2} & \dots & b_{NN} \\
                \end{bmatrix}
            $$

            \noindent The member $c_{ij}$ for $0 \leq i, j \leq N$ can be calculated by the following formula.
            $$c_{ij} = a_{i1} b_{1j} + a_{i2} b_{2j} + \dots + a_{iN} b_{Nj} = \sum_{k = 1}^{N} a_{ik} b_{kj}$$

            \noindent Therefore, matrix multiplication can be implemented as a C code as follows.
            \begin{codeblock}[minted language=c]{Matrix Multiplication}
                #include<stdio.h>
                #define MATRIX_SIZE 8

                int a[MATRIX_SIZE][MATRIX_SIZE];
                int b[MATRIX_SIZE][MATRIX_SIZE];
                int c[MATRIX_SIZE][MATRIX_SIZE];

                int main() {
                    for (int i = 0; i < MATRIX_SIZE; i++) {
                        for (int j = 0; j < MATRIX_SIZE; j++) {
                            c[i][j] = 0;
                            for (int k = 0; k < MATRIX_SIZE; k++) {
                                c[i][j] += a[i][k] * b[k][j];
                            }
                        }
                    }
                    return 0;
                }
            \end{codeblock}

    \pagebreak

        \subsection{No Optimization}
            First, use \texttt{gcc} compiler to compile the C code into assembly by this command line
            \begin{codeblock}[minted language=bat]{Optimization Level 0}
                gcc -O0 -S matrix_multiplication.c -o matrix_multiplication_o0.asm
            \end{codeblock}

            \subsubsection{Stack Memory}
                \begin{itemize}
                    \item Save the base pointer (\texttt{\%rbp}) of the calling function onto the stack.
                    \item Copies the current stack pointer (\texttt{\%rsp}) into the base pointer (\texttt{\%rbp}).
                    \begin{itemize}
                        \item The base pointer (\texttt{\%rbp}) acts as a fixed reference point to access local variables.
                    \end{itemize}
                    \item Subtract 48 from the stack pointer (\texttt{\%rsp}) to allocate 48 bytes of memory to the stack.
                    \begin{itemize}
                        \item 16 bytes memory for \texttt{i, j, k} integers variables.
                        \item 32 bytes memory for function calls.
                    \end{itemize}
                \end{itemize}
                \begin{codeblock}[minted language=asm]{Stack Memory}
                    main:
                        pushq   %rbp
                        movq    %rsp, %rbp
                        subq    $48, %rsp
                \end{codeblock}

                \noindent According to the assembly code above, here is the stack memory partition.
                Note that the \texttt{\%rsp} is the stack pointer (offset = 0)
                and the \texttt{\%rbp} is the base pointer (offset = 48)
                \begin{table}[h]
                    \centering
                    \begin{tabular}{|c|c|l|}
                        \hline
                        Offset from \texttt{\%rsp} & Size (bytes) & Usage \\
                        \hline \hline
                        0 - 31 & 32 & Shadow space for function calls. \\
                        32 - 47 & 16 & Store a local variables (\texttt{i, j, k}) \\
                        \hline
                    \end{tabular}
                \end{table}

    \pagebreak

            \subsubsection{Outer Loop}
                This section is an explanation of the assembly code for the outer loop logic which is this section of the code
                \begin{codeblock}[minted language=c]{Outer Loop}
                    for (int i = 0; i < 8; i++) {
                        for (int j = 0; j < 8; j++) {
                            // MATRIX MULTIPLICATION LOGIC HERE...
                        }
                    }
                    return 0;
                \end{codeblock}

                \begin{itemize}
                    \item Call C runtime initialization function by calling \texttt{\_\_main}
                    \item Store the number \texttt{0} to the \texttt{i} variables which is stored into the
                    memory at 4 bytes before the base pointer (\texttt{\%rbp})
                    \item Then jump to \texttt{.L2}
                \end{itemize}
                \begin{codeblock}[minted language=asm]{Initialize i in the Outer Loop}
                    main:
                        # STACK MEMORY ASM CODE
                        call    __main
                        movl    $0, -4(%rbp)
                        jmp .L2
                \end{codeblock}

                \begin{itemize}
                    \item Check if the \texttt{i} variable is $\leq 7$
                    \begin{itemize}
                        \item \textbf{If yes}, jump to \texttt{.L7}
                        \item \textbf{If no}, continue instructions in \texttt{.L2}
                    \end{itemize}
                    \item Store the number \texttt{0} to the register \texttt{\%eax}
                    \item Add 48 to the stack pointer (\texttt{\%rsp}) to free 48 bytes of memory.
                    \item Restore the base pointer (\texttt{\%rbp}) to the initial value before running this program.
                    \item Ends the program using \texttt{ret} instruction.
                \end{itemize}
                \begin{codeblock}[minted language=asm]{Comparing i in the Outer Loop}
                    .L2:
                        cmpl    $7, -4(%rbp)
                        jle .L7
                        movl    $0, %eax
                        addq    $48, %rsp
                        popq    %rbp
                        ret
                \end{codeblock}

    \pagebreak

                \begin{itemize}
                    \item Store the number \texttt{0} to the \texttt{j} variables which is stored into the
                    memory at 8 bytes before the base pointer (\texttt{\%rbp})
                    \item Then jump to \texttt{.L3}
                \end{itemize}
                \begin{codeblock}[minted language=asm]{Initialize j in the Outer Loop}
                    .L7:
                        movl    $0, -8(%rbp)
                        jmp .L3
                \end{codeblock}

                \begin{itemize}
                    \item Check if the \texttt{j} variable is $\leq 7$
                    \begin{itemize}
                        \item \textbf{If yes}, jump to \texttt{.L6}
                        \item \textbf{If no}, continue instructions in \texttt{.L3}
                    \end{itemize}
                    \item Add the variable \texttt{i} by 1
                \end{itemize}
                \begin{codeblock}[minted language=asm]{Comparing j in the Outer Loop}
                    .L3:
                        cmpl    $7, -8(%rbp)
                        jle .L6
                        addl    $1, -4(%rbp)
                \end{codeblock}

    \pagebreak

            \subsubsection{Inner Loop}
                This section is an explanation of the assembly code for the inner loop logic which is this section of the code
                \begin{codeblock}[minted language=c]{Outer Loop}
                    c[i][j] = 0;
                    for (int k = 0; k < MATRIX_SIZE; k++) {
                        c[i][j] += a[i][k] * b[k][j];
                    }
                \end{codeblock}

                \begin{itemize}
                    \item Handle the variable \texttt{j}
                    \begin{itemize}
                        \item Get the value of the variable \texttt{j} from the memory and put in the register \texttt{\%eax}.
                        \item Extend the value in the register \texttt{\%eax} from 32 bits to 64 bits for memory addressing
                        and store it to the register \texttt{\%rax}.
                    \end{itemize}
                    \item Handle the variable \texttt{i}
                    \begin{itemize}
                        \item Get the value of the variable \texttt{i} from the memory and put in the register \texttt{\%edx}.
                        \item Extend the value of the variable \texttt{i} from the register \texttt{\%edx} from 32 bits to 64 bits
                        for memory addressing and store it to the register \texttt{\%rdx}.
                    \end{itemize}
                    \item Get the memory address of an element in the array.
                    \begin{itemize}
                        \item Calculate $i \times 8$ by shift left the value inside the register \texttt{\%rdx} by 3 bits.
                        \item Add $i \times 8$ from the register \texttt{\%rdx} to the \texttt{j} in the register \texttt{\%rax}.
                        Now, the register \texttt{\%rax} stores $(i \times 8) + j$
                        \item Multiply the value in the register \texttt{\%rax} by 4 because an integer takes 4 bytes of memory.
                        Then, store the result $4((i \times 8) + j)$ to the register \texttt{\%rdx}
                        \item Store the starting memory address of matrix \texttt{c} to the register \texttt{\%rax}
                        \item Store the number 0 to the \texttt{c[i][j]} where memory address is at the sum of the
                        starting memory address (\texttt{\%rax}) and the offset (\texttt{\%rdx}).
                    \end{itemize}
                \end{itemize}
                \begin{codeblock}[minted language=asm]{Initialize \texttt{c[i][j]}}
                    .L6:
                        movl    -8(%rbp), %eax
                        cltq
                        movl    -4(%rbp), %edx
                        movslq  %edx, %rdx
                        salq    $3, %rdx
                        addq    %rdx, %rax
                        leaq    0(,%rax,4), %rdx
                        leaq    c(%rip), %rax
                        movl    $0, (%rdx,%rax)
                \end{codeblock}

    \pagebreak

                \begin{itemize}
                    \item Store the number \texttt{0} to the \texttt{k} variables which is stored into the
                    memory at 12 bytes before the base pointer (\texttt{\%rbp})
                    \item Then jump to \texttt{.L4}
                    \item Check if the \texttt{k} variable is $\leq 7$
                    \begin{itemize}
                        \item \textbf{If yes}, jump to \texttt{.L5}
                        \item \textbf{If no}, continue instructions in \texttt{.L4}
                    \end{itemize}
                    \item Add the variable \texttt{j} by 1
                \end{itemize}
                \begin{codeblock}[minted language=asm]{Inner Loop}
                    .L6:
                        # Initialize c[i][j] = 0 instruction in ASM
                        movl    $0, -12(%rbp)
                        jmp .L4

                    .L4:
                        cmpl    $7, -12(%rbp)
                        jle .L5
                        addl    $1, -8(%rbp)
                \end{codeblock}

                Then, access the value of \texttt{c[i][j]} by using the similar steps from the page 5,
                and store the value of \texttt{c[i][j]} to the register \texttt{\%ecx}
                \begin{codeblock}[minted language=asm]{Access \texttt{c[i][j]}}
                    .L5:
                        # Store i and j to the registers
                        movl    -8(%rbp), %eax
                        cltq
                        movl    -4(%rbp), %edx
                        movslq  %edx, %rdx
                        # Calculate the memory address of c[i][j]
                        salq    $3, %rdx
                        addq    %rdx, %rax
                        leaq    0(,%rax,4), %rdx
                        leaq    c(%rip), %rax
                        # Store the value of c[i][j] to the register
                        movl    (%rdx,%rax), %ecx
                \end{codeblock}

    \pagebreak

                Next, access the value of \texttt{a[i][k]} by using the similar steps from the page 5,
                and store the value of \texttt{a[i][k]} to the register \texttt{\%edx}
                \begin{codeblock}[minted language=asm]{Access \texttt{a[i][k]}}
                    .L5:
                        # PREVIOUS INSTRUCTIONS...

                        # Store i and k to the registers
                        movl    -12(%rbp), %eax
                        cltq
                        movl    -4(%rbp), %edx
                        movslq  %edx, %rdx
                        # Calculate the memory address of a[i][k]
                        salq    $3, %rdx
                        addq    %rdx, %rax
                        leaq    0(,%rax,4), %rdx
                        leaq    a(%rip), %rax
                        # Store the value of a[i][k] to the register
                        movl    (%rdx,%rax), %edx
                \end{codeblock}

                After accessing \texttt{a[i][k]}, access the value of \texttt{b[k][j]}  by using the similar steps
                from the page 5, and store the value of \texttt{b[k][j]} to the register \texttt{\%eax}
                \begin{codeblock}[minted language=asm]{Access \texttt{b[k][j]}}
                    .L5:
                        # PREVIOUS INSTRUCTIONS...

                        # Store j and k to the registers
                        movl    -8(%rbp), %eax
                        cltq
                        movl    -12(%rbp), %r8d
                        movslq  %r8d, %r8
                        # Calculate the memory address of b[k][j]
                        salq    $3, %r8
                        addq    %r8, %rax
                        leaq    0(,%rax,4), %r8
                        leaq    b(%rip), %rax
                        # Store the value of b[k][j] to the register
                        movl    (%r8,%rax), %eax
                \end{codeblock}

    \pagebreak

            \subsubsection{Matrix Multiplication}
                This section is an explanation of the assembly code for the matrix multiplication logic which is this section of the code
                \begin{codeblock}[minted language=c]{Matrix Multiplication}
                    c[i][j] += a[i][k] * b[k][j];
                \end{codeblock}

                Before explaining the logic, remind that we have stores these values inside the registers
                in the previous section
                \begin{itemize}
                    \item The value of \texttt{c[i][j]} is stored in the register \texttt{\%ecx}
                    \item The value of \texttt{a[i][k]} is stored in the register \texttt{\%edx}
                    \item The value of \texttt{b[k][j]} is stored in the register \texttt{\%eax}
                \end{itemize}

                \noindent This is the further explanation of assembly code of \texttt{.L5}
                \begin{itemize}
                    \item Mutiply \texttt{a[i][k]} from the register \texttt{\%edx} with
                    \texttt{b[k][j]} from the register \texttt{\%eax}.
                    Then, store the product to the register \texttt{\%eax}
                    \item Add the product \texttt{a[i][k] * b[k][j]} from the register \texttt{\%eax}
                    to \texttt{c[i][j]} from the register \texttt{\%ecx}.
                    Then, store the sum to the register \texttt{\%ecx}.
                \end{itemize}
                \begin{codeblock}[minted language=asm]{Math Process}
                    .L5:
                        # PREVIOUS INSTRUCTIONS...
                        imull   %edx, %eax
                        addl    %eax, %ecx
                \end{codeblock}

                Finally, calculate the memory address of \texttt{c[i][j]} and
                store the result \texttt{c[i][j] + (a[i][k] * b[k][j])} which is calculated in the previous step
                to the \texttt{c[i][j]} and increment \texttt{k} by 1.
                \begin{codeblock}[minted language=asm]{Store the Result}
                    .L5:
                        # Store i and j to the registers
                        movl    -8(%rbp), %eax
                        cltq
                        movl    -4(%rbp), %edx
                        movslq  %edx, %rdx
                        # Calculate the memory address of c[i][j]
                        salq    $3, %rdx
                        addq    %rdx, %rax
                        leaq    0(,%rax,4), %rdx
                        leaq    c(%rip), %rax
                        # Store the result to c[i][j]
                        movl    %ecx, (%rdx,%rax)
                        # Increment k by 1
                        addl    $1, -12(%rbp)
                \end{codeblock}

    \pagebreak

        \subsection{With Optimization}
            First, use \texttt{gcc} compiler to compile the C code into assembly by this command line
            \begin{codeblock}[minted language=bat]{Optimization Level 3}
                gcc -O3 -S matrix_multiplication.c -o matrix_multiplication_o3.asm
            \end{codeblock}

            \subsubsection{Stack Memory}
                First, subtract 248 from the stack pointer \texttt{\%rsp} to allocate 248 bytes of memory to the stack.
                Then, partition the memory for backup the state of vector registers \texttt{\%xxm6 - \%xxm15}.
                \begin{table}[h]
                    \centering
                    \begin{tabular}{|c|c|l|}
                        \hline
                        Offset from \texttt{\%rsp} & Size (bytes) & Usage \\
                        \hline \hline
                        0 - 31 & 32 & Shadow space for function calls. \\
                        32 - 79 & 48 & Store a local variables or temporary math results \\
                        80 - 95 & 16 & Save the state of vector register \texttt{\%xx6} \\
                        96 - 111 & 16 & Save the state of vector register \texttt{\%xx7} \\
                        112 - 127 & 16 & Save the state of vector register \texttt{\%xx8} \\
                        128 - 143 & 16 & Save the state of vector register \texttt{\%xx9} \\
                        144 - 159 & 16 & Save the state of vector register \texttt{\%xx10} \\
                        160 - 175 & 16 & Save the state of vector register \texttt{\%xx11} \\
                        176 - 191 & 16 & Save the state of vector register \texttt{\%xx12} \\
                        192 - 207 & 16 & Save the state of vector register \texttt{\%xx13} \\
                        208 - 223 & 16 & Save the state of vector register \texttt{\%xx14} \\
                        224 - 239 & 16 & Save the state of vector register \texttt{\%xx15} \\
                        240 - 247 & 8 & 16 bytes alignment padding \\
                        \hline
                    \end{tabular}
                \end{table}
                \begin{codeblock}[minted language=asm]{Stack Memory Partition}
                    main:
                        subq    $248, %rsp
                        movups  %xmm6, 80(%rsp)
                        movups  %xmm7, 96(%rsp)
                        movups  %xmm8, 112(%rsp)
                        movups  %xmm9, 128(%rsp)
                        movups  %xmm10, 144(%rsp)
                        movups  %xmm11, 160(%rsp)
                        movups  %xmm12, 176(%rsp)
                        movups  %xmm13, 192(%rsp)
                        movups  %xmm14, 208(%rsp)
                        movups  %xmm15, 224(%rsp)
                \end{codeblock}

    \pagebreak
            \subsubsection{Pointers}
                Next, call C runtime initialization function by calling \texttt{\_\_main}
                and store the pointers to the registers.
                \begin{table}[h]
                    \centering
                    \begin{tabular}{|c|l|}
                        \hline
                        Register & Usage \\
                        \hline \hline
                        \texttt{\%rdx} & The starting memory address of matrix \texttt{a} \\
                        \texttt{\%rcx} & The starting memory address of matrix \texttt{c} \\
                        \texttt{\%r8} & The ending memory address of matrix \texttt{a} \\
                        \texttt{\%rax} & The starting memory address of matrix \texttt{b} \\
                        \hline
                    \end{tabular}
                \end{table}

                Note that the matrix has $8 \times 8$ which has 64 integers and takes 256 bytes of memory.
                Therefore, adding 256 to starting memory address of matrix \texttt{a} will get
                the ending memory address of matrix \texttt{a}.

                \begin{codeblock}[minted language=asm]{Pointers}
                    main:
                        call    __main
                        leaq    a(%rip), %rdx
                        leaq    c(%rip), %rcx
                        leaq    256(%rdx), %r8
                        leaq    b(%rip), %rax
                \end{codeblock}

    \pagebreak

            \subsubsection{Loop Unrolling \& Instruction Scheduling}
                In level 3 optimization of GCC compiler, it performing loop unrolling.
                which can be seen in the assembly code that it is hardcoded and
                has no loop at all.

                \begin{codeblock}[minted language=asm]{Matrix Calculation Instructions}
                    .L2:
                        movd    (%rdx), %xmm3
                        movdqu  (%rax), %xmm1
                        addq    $32, %rdx
                        addq    $32, %rcx
                        movd    -28(%rdx), %xmm4
                        movd    -24(%rdx), %xmm5
                        pshufd  $0, %xmm3, %xmm8
                        movdqa  %xmm8, %xmm13
                        psrlq   $32, %xmm1
                        movdqa  %xmm8, %xmm9
                        # MORE ASM INSTRUCTIONS...
                \end{codeblock}

                Moreover, the instructions are out of order compared to the original C code
                and seems random because GCC compiler is performing the instruction scheduling
                which reduces total clock cycle.

                Since loading data takes more clock cycle than the mathematics operations,
                the GCC compiler fill this gap by scheduling instructions which is unrelated
                to the current calculation. This is an example from the actual assembly code.
                \begin{itemize}
                    \item Load 4 integers into the vector register \texttt{\%xmm1}.
                    \item These are instructions which is unrelated to the vector register \texttt{\%xmm1}
                    \begin{itemize}
                        \item Add 32 to the pointer \texttt{\%rdx} and \texttt{\%rcx}.
                        \item Load an integer to vector register \texttt{\%xmm4} then
                        load another integer to vector register \texttt{\%xmm5}.
                        \item Fill all slots inside the vector register \texttt{\%xmm8} with
                        the first element from the vector register \texttt{\%xmm3}.
                        \item Copy all elements from the vector register \texttt{\%xmm8} to
                        the vector register \texttt{\%xmm13}.
                    \end{itemize}
                    \item Finally performing an operations to the vector register \texttt{\%xmm1}
                    which is shift right by 32 bits.
                \end{itemize}
                \begin{codeblock}[minted language=asm]{Instruction Order}
                    .L2:
                        movdqu  (%rax), %xmm1
                        addq    $32, %rdx
                        addq    $32, %rcx
                        movd    -28(%rdx), %xmm4
                        movd    -24(%rdx), %xmm5
                        pshufd  $0, %xmm3, %xmm8
                        movdqa  %xmm8, %xmm13
                        psrlq   $32, %xmm1
                        # MORE ASM INSTRUCTIONS...
                \end{codeblock}

    \pagebreak

            \subsubsection{Vectorization}
                In addition, the \texttt{-O3} optimization also performs vectorization by using several SIMD instructions.
                \begin{itemize}
                    \item Data movement \& manipulation instructions
                    \begin{itemize}
                        \item The vector registers \texttt{\%xmm6} to \texttt{\%xmm15} have 128 bits of memory
                        which contains 4 integers at once
                        \item The \texttt{movdqu} and \texttt{movdqa} instruction loads 4 integers to the vector register at once.
                        \item The \texttt{pushfd} instruction copy an integer from a vector register and
                        fill that integer to all slots in another vector register at once.
                        \item The \texttt{prslq} shift right bits on all value in register at once.
                        \item The \texttt{punpckldq} merges two vector registers.
                    \end{itemize}
                    \item Mathematics operation instructions
                    \begin{itemize}
                        \item The \texttt{pmuludq} instruction multiples two vector registers at once.
                        \item The \texttt{paddd} instruction adds 4 integers inside the vector register at once.
                    \end{itemize}
                \end{itemize}
                \begin{codeblock}[minted language=asm]{SIMD Instruction}
                    .L2:
                        movdqu  (%rax), %xmm1
                        # OTHER ASM INSTRUCTIONS...
                        pshufd  $0, %xmm3, %xmm8
                        movdqa  %xmm8, %xmm13
                        psrlq   $32, %xmm1
                        # OTHER ASM INSTRUCTIONS...
                        pmuludq b(%rip), %xmm9
                        # OTHER ASM INSTRUCTIONS...
                        punpckldq    %xmm1, %xmm9
                        # OTHER ASM INSTRUCTIONS...
                        paddd    %xmm10, %xmm9
                        # OTHER ASM INSTRUCTIONS...
                \end{codeblock}

    \pagebreak

            \subsection{Summary}

                \noindent \textbf{No Optimization (\texttt{gcc -O0})}
                \begin{itemize}
                    \item \textbf{Basic Blocks}: The assembly code is separated into several basic blocks
                    (e.g. \texttt{.L3, .L4, .L5}) which are connected by jump instruction (\texttt{jmp, jle}).
                    \item \textbf{High Control Overhead}: For every single mathematical calculation, the processor
                    must update the counter \texttt{i, j, k} and compare if it is $\leq 7$.
                \end{itemize}

                \noindent \textbf{With Optimization (\texttt{gcc -O3})}
                \begin{itemize}
                    \item \textbf{Loop Unrolling}: All matrix calculation loops are replaced by a massive instructions
                    inside \texttt{.L2} block instead of jumping between blocks like the \texttt{gcc -O0} version.
                    \item \textbf{Vectorization}: It uses vector registers (\texttt{\%xmm6 - \%xmm15}) for loading
                    data and do the mathematical operations with SIMD instructions.
                    \item \textbf{Loop Interchange}: The compiler performed loop interchange to improve locality,
                    ensuring the data could be smoothly fed into those vector registers and making vectorization work efficiently
                \end{itemize}

    \pagebreak

    \section{Problem II}
        Try this simple code. See if the compiler is able to perform any optimization.
        Please provide your analysis.
        \begin{codeblock}[minted language=c]{Simple Code}
            #include<stdio.h>
            #define MATRIX_SIZE 8

            double a[MATRIX_SIZE][MATRIX_SIZE];
            double b[MATRIX_SIZE][MATRIX_SIZE];

            int main() {
                for (int i = 0; i < MATRIX_SIZE; i++) {
                    for (int j = 0; j < MATRIX_SIZE; j++) {
                        a[i][j] = b[i][j] + 2.5;
                    }
                }
                return 0;
            }
        \end{codeblock}

	\pagebreak

        \subsection{No Optimization}
            First, use \texttt{gcc} compiler to compile the C code into assembly by this command line
            \begin{codeblock}[minted language=bat]{Optimization Level 0}
                gcc -O0 -S simple_code.c -o simple_code_o0.asm
            \end{codeblock}

            \subsubsection{Stack Memory}
                \begin{itemize}
                    \item Save the base pointer (\texttt{\%rbp}) of the calling function onto the stack.
                    \item Copies the current stack pointer (\texttt{\%rsp}) into the base pointer (\texttt{\%rbp}).
                    \begin{itemize}
                        \item The base pointer (\texttt{\%rbp}) acts as a fixed reference point to access local variables.
                    \end{itemize}
                    \item Subtract 48 from the stack pointer (\texttt{\%rsp}) to allocate 48 bytes of memory to the stack.
                    \begin{itemize}
                        \item 16 bytes memory for \texttt{i, j} integers variables.
                        \item 32 bytes memory for function calls.
                    \end{itemize}
                \end{itemize}
                \begin{codeblock}[minted language=asm]{Stack Memory}
                    main:
                        pushq   %rbp
                        movq    %rsp, %rbp
                        subq    $48, %rsp
                        call    __main
                        movl    $0, -4(%rbp)
                        jmp	.L2
                \end{codeblock}

                \noindent According to the assembly code above, here is the stack memory partition.
                Note that the \texttt{\%rsp} is the stack pointer (offset = 0)
                and the \texttt{\%rbp} is the base pointer (offset = 48)
                \begin{table}[h]
                    \centering
                    \begin{tabular}{|c|c|l|}
                        \hline
                        Offset from \texttt{\%rsp} & Size (bytes) & Usage \\
                        \hline \hline
                        0 - 31 & 32 & Shadow space for function calls. \\
                        32 - 47 & 16 & Store a local variables (\texttt{i, j}) \\
                        \hline
                    \end{tabular}
                \end{table}

    \pagebreak

            \subsubsection{Nested Loop}
                This section is an explanation of the assembly code for the outer loop logic which is this section of the code
                \begin{codeblock}[minted language=c]{Nested Loop}
                    for (int i = 0; i < MATRIX_SIZE; i++) {
                        for (int j = 0; j < MATRIX_SIZE; j++) {
                            # MATRIX CALCULATION LOGIC HERE...
                        }
                    }
                    return 0;
                \end{codeblock}

                \begin{itemize}
                    \item Call C runtime initialization function by calling \texttt{\_\_main}
                    \item Store the number \texttt{0} to the \texttt{i} variables which is stored into the
                    memory at 4 bytes before the base pointer (\texttt{\%rbp})
                    \item Then jump to \texttt{.L2}
                \end{itemize}
                \begin{codeblock}[minted language=asm]{Initialize i in the Outer Loop}
                    main:
                        # STACK MEMORY ASM CODE
                        call    __main
                        movl    $0, -4(%rbp)
                        jmp .L2
                \end{codeblock}

                \begin{itemize}
                    \item Check if the \texttt{i} variable is $\leq 7$
                    \begin{itemize}
                        \item \textbf{If yes}, jump to \texttt{.L5}
                        \item \textbf{If no}, continue instructions in \texttt{.L2}
                    \end{itemize}
                    \item Store the number \texttt{0} to the register \texttt{\%eax}
                    \item Add 48 to the stack pointer (\texttt{\%rsp}) to free 48 bytes of memory.
                    \item Restore the base pointer (\texttt{\%rbp}) to the initial value before running this program.
                    \item Ends the program using \texttt{ret} instruction.
                \end{itemize}
                \begin{codeblock}[minted language=asm]{Comparing i in the Outer Loop}
                    .L2:
                        cmpl    $7, -4(%rbp)
                        jle .L5
                        movl    $0, %eax
                        addq    $48, %rsp
                        popq    %rbp
                        ret
                \end{codeblock}

    \pagebreak

                \begin{itemize}
                    \item Store the number \texttt{0} to the \texttt{j} variables which is stored into the
                    memory at 8 bytes before the base pointer (\texttt{\%rbp})
                    \item Then jump to \texttt{.L3}
                \end{itemize}
                \begin{codeblock}[minted language=asm]{Initialize j in the Inner Loop}
                    .L5:
                        movl    $0, -8(%rbp)
                        jmp .L3
                \end{codeblock}

                \begin{itemize}
                    \item Check if the \texttt{j} variable is $\leq 7$
                    \begin{itemize}
                        \item \textbf{If yes}, jump to \texttt{.L4}
                        \item \textbf{If no}, continue instructions in \texttt{.L3}
                    \end{itemize}
                    \item Add the variable \texttt{i} by 1
                \end{itemize}
                \begin{codeblock}[minted language=asm]{Comparing j in the Inner Loop}
                    .L3:
                        cmpl    $7, -8(%rbp)
                        jle .L4
                        addl    $1, -4(%rbp)
                \end{codeblock}

    \pagebreak

            \subsubsection{Matrix Addition}
                This section is an explanation of the assembly code for the matrix addition logic which is this section of the code
                \begin{codeblock}[minted language=c]{Matrix Addition}
                    a[i][j] = b[i][j] + 2.5;
                \end{codeblock}

                \begin{itemize}
                    \item Handle the variable \texttt{j}
                    \begin{itemize}
                        \item Get the value of the variable \texttt{j} from the memory and put in the register \texttt{\%eax}.
                        \item Extend the value in the register \texttt{\%eax} from 32 bits to 64 bits for memory addressing
                        and store it to the register \texttt{\%rax}.
                    \end{itemize}
                    \item Handle the variable \texttt{i}
                    \begin{itemize}
                        \item Get the value of the variable \texttt{i} from the memory and put in the register \texttt{\%edx}.
                        \item Extend the value of the variable \texttt{i} from the register \texttt{\%edx} from 32 bits to 64 bits
                        for memory addressing and store it to the register \texttt{\%rdx}.
                    \end{itemize}
                    \item Get the memory address of an element in the array.
                    \begin{itemize}
                        \item Calculate $i \times 8$ by shift left the value inside the register \texttt{\%rdx} by 3 bits.
                        \item Add $i \times 8$ from the register \texttt{\%rdx} to the \texttt{j} in the register \texttt{\%rax}.
                        Now, the register \texttt{\%rax} stores $(i \times 8) + j$
                        \item Multiply the value in the register \texttt{\%rax} by 8 because a single \texttt{double} value takes 8 bytes of memory.
                        Then, store the result $8((i \times 8) + j)$ to the register \texttt{\%rdx}
                        \item Store the starting memory address of matrix \texttt{b} to the register \texttt{\%rax}
                        \item Store the memory address to the \texttt{b[i][j]} where memory address is at the sum of the
                        starting memory address (\texttt{\%rax}) and the offset (\texttt{\%rdx}) inside the register \texttt{\%xmm1}.
                    \end{itemize}
                \end{itemize}
                \begin{codeblock}[minted language=asm]{Access Memory Address}
                    .L4:
                        movl    -8(%rbp), %eax
                        cltq
                        movl    -4(%rbp), %edx
                        movslq  %edx, %rdx
                        salq    $3, %rdx
                        addq    %rdx, %rax
                        leaq    0(,%rax,8), %rdx
                        leaq    b(%rip), %rax
                        movsd   (%rdx,%rax), %xmm1
                \end{codeblock}

    \pagebreak

                \begin{itemize}
                    \item Store a double-precision constant number 2.5 from \texttt{.LC0} to the register {\%xmm0}.
                    \item Add the \texttt{b[i][j]} from the register \texttt{\%xmm1} to the register \texttt{\%xmm0}.
                    Now, the register \texttt{\%xmm0} contains \texttt{b[i][j] + 2.5}
                \end{itemize}
                \begin{codeblock}[minted language=asm]{Matrix Addition}
                    .L4:
                        # PREVIOUS ASM OPERATIONS...
                        movsd   .LC0(%rip), %xmm0
                        addsd   %xmm1, %xmm0

                    .LC0:
                        .long   0
                        .long   1074003968
                        .def    __main
                        .ident  "GCC: (Rev7, Built by MSYS2 project) 15.1.0"
                \end{codeblock}

                Finally, calculate the memory address of \texttt{a[i][j]} and
                store the result \texttt{b[i][j] + 2.5} which is in the register
                \texttt{\%xmm0} and increment \texttt{j} by 1.
                \begin{codeblock}[minted language=asm]{Store the Result}
                    .L4:
                        # PREVIOUS ASM OPERATIONS...

                        # Store i and j to the registers
                        movl    -8(%rbp), %eax
                        cltq
                        movl    -4(%rbp), %edx
                        movslq  %edx, %rdx
                        # Calculate the memory address of a[i][j]
                        salq    $3, %rdx
                        addq    %rdx, %rax
                        leaq    0(,%rax,8), %rdx
                        leaq    a(%rip), %rax
                        # Store the result to a[i][j]
                        movsd   %xmm0, (%rdx,%rax)
                        # Increment j by 1
                        addl    $1, -8(%rbp)
                \end{codeblock}

    \pagebreak

        \subsection{With Optimization}
            First, use \texttt{gcc} compiler to compile the C code into assembly by this command line
            \begin{codeblock}[minted language=bat]{Optimization Level 3}
                gcc -O3 -S simple_code.c -o simple_code_o3.asm
            \end{codeblock}

            In level 3 optimization of GCC compiler, it performing loop unrolling.
            which can be seen in the assembly code that it is hardcoded and
            has no loop at all.

            \vspace{1em}

            This assembly code also uses SIMD instruction to perform vectorization.
            According to the assembly code, it do as follows:
            \begin{itemize}
                \item It uses vector registers \texttt{\%xmm0 - \%xmm4} to store multiple numbers at once.
                Each has 128 bits memory which can store 2 double numbers.
                \item Store a double-precision constant number 2.5 from \texttt{.LC1} to the register {\%xmm0}.
                \item SIMD data loading instruction - \texttt{movupd}
                \begin{itemize}
                    \item Load \texttt{b[0][2]} and \texttt{b[0][3]} to the vector register \texttt{\%xmm3} at once.
                    \item Load \texttt{b[0][4]} and \texttt{b[0][5]} to the vector register \texttt{\%xmm2} at once.
                    \item Load \texttt{b[0][6]} and \texttt{b[0][7]} to the vector register \texttt{\%xmm1} at once.
                    \item Load \texttt{b[0][0]} and \texttt{b[0][1]} to the vector register \texttt{\%xmm4} at once.
                \end{itemize}

                \item SIMD adding instruction - \texttt{addpd}
                \begin{itemize}
                    \item Add \texttt{b[0][2]} and \texttt{b[0][3]} with 2.5 at once.
                    \item Add \texttt{b[0][4]} and \texttt{b[0][5]} with 2.5 at once.
                    \item Add \texttt{b[0][6]} and \texttt{b[0][7]} with 2.5 at once.
                    \item Add \texttt{b[0][0]} and \texttt{b[0][1]} with 2.5 at once.
                \end{itemize}

                \item SIMD data storing instruction - \texttt{movups}
                \begin{itemize}
                    \item Store \texttt{b[0][2] + 2.5} and \texttt{b[0][3] + 2.5} back to the matrix \texttt{b} at once.
                    \item Store \texttt{b[0][4] + 2.5} and \texttt{b[0][5] + 2.5} back to the matrix \texttt{b} at once.
                    \item Store \texttt{b[0][0] + 2.5} and \texttt{b[0][1] + 2.5} back to the matrix \texttt{b} at once.
                    \item Store \texttt{b[0][6] + 2.5} and \texttt{b[0][7] + 2.5} back to the matrix \texttt{b} at once.
                \end{itemize}
            \end{itemize}

            The assembly code is in the next page $\rightarrow$

    \pagebreak

            \begin{codeblock}[minted language=asm]{Matrix Addition Instructions}
                main:
                    # Load a constant number 2.5
                    movddup .LC1(%rip), %xmm0
                    # Load numbers to the register
                    movupd  16+b(%rip), %xmm3
                    movupd  32+b(%rip), %xmm2
                    movupd  48+b(%rip), %xmm1
                    movupd  b(%rip), %xmm4
                    # Add 2.5 to the register
                    addpd   %xmm0, %xmm3
                    addpd   %xmm0, %xmm2
                    addpd   %xmm0, %xmm1
                    addpd   %xmm0, %xmm4
                    # Store back to the memory
                    movups  %xmm3, 16+a(%rip)
                    movups  %xmm2, 32+a(%rip)
                    movups  %xmm4, a(%rip)
                    movups  %xmm1, 48+a(%rip)

                    # MORE ASM INSTRUCTIONS...
                
                .LC1:
                    .long   0
                    .long   1074003968
                    .def    __main
                    .ident  "GCC: (Rev7, Built by MSYS2 project) 15.1.0"
            \end{codeblock}

    \pagebreak

            \subsection{Summary}

                \noindent \textbf{No Optimization (\texttt{gcc -O0})}
                \begin{itemize}
                    \item \textbf{Basic Blocks}: The assembly code is separated into several basic blocks
                    (e.g. \texttt{.L2, .L3}) which are connected by jump instruction (\texttt{jmp, jle}).
                    \item \textbf{High Control Overhead}: For every single mathematical calculation, the processor
                    must update the counter \texttt{i, j} and compare if it is $\leq 7$.
                \end{itemize}

                \noindent \textbf{With Optimization (\texttt{gcc -O3})}
                \begin{itemize}
                    \item \textbf{Loop Unrolling}: All matrix calculation loops are replaced by a massive instructions
                    inside \texttt{main} block instead of jumping between blocks like the \texttt{gcc -O0} version.
                    \item \textbf{Vectorization}: It uses vector registers (\texttt{\%xmm0 - \%xmm4}) for loading
                    data and do the mathematical operations with SIMD instructions.
                \end{itemize}

\end{document}